\begin{table}[H]
  \centering
  \caption{Alcances y limitaciones del estudio}
  \label{tab:alcances}
  \begin{tabular}{@{}p{0.42\textwidth}p{0.48\textwidth}@{}}
    \toprule
    \textbf{Alcance} & \textbf{Limitación} \\
    \midrule
    Evaluar alineación \textbf{estructural} y \textbf{semántica} entre contratos OpenAPI y modelos BIAN &
    No desarrollar ni desplegar APIs en entornos productivos \\
    Definir \textbf{métricas reproducibles} aplicables sobre artefactos documentales &
    No sustituir procesos de certificación regulatoria ni auditoría legal \\
    Definir un \textbf{modelo de alineación} (S, E, C $\rightarrow$ \AlignmentScore{}) aplicable a contratos OpenAPI bancarios &
    No evaluar el catálogo completo de \ServiceDomain{}s BIAN \\
    Analizar contratos OpenAPI del sector bancario frente a \ServiceDomain{}s BIAN de referencia &
    No realizar encuestas ni estudios de percepción a desarrolladores \\
    Integrar hallazgos en \AlignmentScore{} y clasificación por niveles de alineación &
    No abordar como foco principal la calidad de testing de APIs (Ehsan et al., 2022; Golmohammadi et al., 2024) \\
    Documentar el \textbf{diseño} del modelo de alineación y del procedimiento de cálculo &
    No incluir la \textbf{construcción} del artefacto software ni la \textbf{experimentación} empírica a escala (fuera del alcance de este plan) \\
    \bottomrule
  \end{tabular}
\end{table}
